\section{Cinemática Inversa}
\label{sec:inversa}


A cinemática inversa faz o caminho oposto ao da cinemática direta: dada a posição pretendida para o centro de preensão da garra, no referencial da base, determina os ângulos das juntas que a produzem. É este o problema que o robô tem de resolver sempre que queremos que a garra vá a um ponto do espaço. Ao contrário da cinemática direta, pode não ter solução, quando o ponto está fora do alcance do braço, ou ter várias. No nosso caso, resolvemo-la de forma analítica, aproveitando as expressões explícitas da secção~\ref{sec:explicita}, o que dá a solução diretamente, sem iterações.

%---------------------------------------------------------------------
\subsection{Formulação do problema}

A função \texttt{inverse\_kin.m}, do pacote \texttt{+movement}, recebe a posição pretendida $\mathbf{p}=(x,y,z)$ do centro de preensão, em metros e no referencial da base, o vetor \texttt{Links} e, opcionalmente, os limites das juntas, a rotação de um objeto, a inclinação preferida da garra e a família de soluções. Devolve as seis juntas $\mathbf{q}=(q_1,\dots,q_6)$ em radianos, um valor lógico que indica se o ponto é alcançável e a inclinação da garra que foi efetivamente usada, na forma \texttt{[armJoints, isReachable, pitchUsed] = inverse\_kin(x, y, z, Links, \dots)}.

A posição impõe três equações, mas as juntas J1 a J4 são quatro incógnitas, pelo que sobra um grau de liberdade. Resolvemo-lo escolhendo a inclinação da garra $\psi=-\varphi$, definida na secção~\ref{sec:explicita}, cujo valor preferido é $\psi=90^\circ$, isto é, a garra na vertical a apontar para baixo, a melhor orientação para apanhar objetos. A junta J5 fica sempre a $0$, para a garra se manter no plano do braço, e a junta J6 apenas orienta a garra em torno do seu eixo, sem alterar a posição do centro de preensão.

%---------------------------------------------------------------------
\subsection{Resolução analítica}

\paragraph{Junta J1.} Pelas expressões da cinemática direta, o ponto $(x,y)$ é o vetor $(L_5,h)$ rodado de $q_1$ em torno do eixo $z$:
\begin{align*}
x &= L_5\cos q_1 - h\sin q_1,\\
y &= L_5\sin q_1 + h\cos q_1.
\end{align*}
Uma rotação conserva o comprimento do vetor, logo $x^2+y^2=L_5^2+h^2$. Com $r=\sqrt{x^2+y^2}$, a distância em frente, no plano do braço, é
\begin{equation*}
h=\sqrt{r^2-L_5^2},
\end{equation*}
que só existe se $r>L_5$. A rotação soma o ângulo $q_1$ ao ângulo do vetor $(L_5,h)$, pelo que $\operatorname{atan2}(y,x)=q_1+\operatorname{atan2}(h,L_5)$ e
\begin{equation*}
q_1=\operatorname{atan2}(y,x)-\operatorname{atan2}(h,L_5).
\end{equation*}
É por causa do desvio lateral $L_5$ que o segundo termo aparece: sem ele, o plano do braço passaria pelo eixo de J1 e bastaria $q_1=\operatorname{atan2}(y,x)$. Tomámos a raiz positiva para $h$, que corresponde ao braço virado para o ponto, em frente da base.

\paragraph{Ponto do pulso.} Fixada a inclinação, $\varphi=-\psi$, e conhecido o ponto da garra no plano do braço, $(h,\,z-L_1)$, retiramos-lhe o elo J4\,$\to$\,J5 e a distância $D=L_6+L_7$, que dependem apenas de $\varphi$. Ficamos com o ponto do pulso (J4), medido a partir de J2:
\begin{align*}
W_h &= h + L_4\sin\varphi - D\cos\varphi,\\
W_v &= (z-L_1) - L_4\cos\varphi - D\sin\varphi.
\end{align*}

\paragraph{Juntas J2 e J3.} Entre J2 e o pulso resta um braço planar de dois elos, de comprimentos $L_2$ e $L_3$. Chamando $\alpha_1=q_2+\pi/2$ ao ângulo do braço superior medido a partir do eixo horizontal, as expressões da cinemática direta escrevem-se
\begin{align*}
W_h &= L_2\cos\alpha_1 + L_3\cos(\alpha_1+q_3),\\
W_v &= L_2\sin\alpha_1 + L_3\sin(\alpha_1+q_3).
\end{align*}
Somando os quadrados, obtém-se a lei dos cossenos no triângulo J2--J3--pulso:
\begin{equation*}
\cos q_3=\frac{W_h^2+W_v^2-L_2^2-L_3^2}{2\,L_2\,L_3}.
\end{equation*}
Para existir solução é preciso $|\cos q_3|\le 1$, o que equivale a pedir que a distância do pulso a J2 esteja entre $|L_2-L_3|$ e $L_2+L_3$. No nosso robô, isso dá entre $0{,}0144~\mathrm{m}$ e $0{,}2064~\mathrm{m}$. Resolvendo em ordem a $q_3$ e depois a $\alpha_1$, vem
\begin{align*}
q_3 &= \pm\arccos\!\left(\cos q_3\right),\\
\alpha_1 &= \operatorname{atan2}(W_v,W_h)-\operatorname{atan2}\!\left(L_3\sin q_3,\;L_2+L_3\cos q_3\right),\\
q_2 &= \alpha_1-\frac{\pi}{2}.
\end{align*}

\paragraph{Junta J4 e juntas J5 e J6.} Como $\varphi=q_2+q_3+q_4$, a última junta do braço planar é
\begin{equation*}
q_4=\varphi-q_2-q_3.
\end{equation*}
Por fim, $q_5=0$. Para $q_6$, quando conhecemos a rotação $\theta_{obj}$ de um cubo em torno de $z$, rodamos a garra de modo a que os dedos fechem sobre duas faces opostas. Como os dedos fecham segundo a direção definida por J1 e J6, e um cubo se repete de $90^\circ$ em $90^\circ$, escolhemos a menor rotação de J6:
\begin{equation*}
q_6 = \left[\left(q_1 - \theta_{\text{obj}} + \frac{\pi}{4}\right) \bmod \frac{\pi}{2}\right] - \frac{\pi}{4}\end{equation*}
que fica entre $-45^\circ$ e $45^\circ$. Só fazemos isto se a garra estiver quase vertical ($\psi>45^\circ$); caso contrário, $q_6=0$.

%---------------------------------------------------------------------
\subsection{Soluções, limites e escolha da inclinação}

\paragraph{Duas soluções para J3.} O sinal de $\pm\arccos$ dá duas soluções para o mesmo ponto: o cotovelo para cima ($q_3\le 0$) e o cotovelo para baixo ($q_3\ge 0$). No nosso caso, fizemos desta maneira: testamos as duas e, entre as válidas, ficamos com a de cotovelo mais alto, isto é, a de maior altura $z_{c}=L_1+L_2\cos q_2$, que mantém o braço afastado da mesa. O código prevê ainda uma família «atrás», com o braço inclinado para trás por cima da base, que não usamos por defeito.

\paragraph{Soluções rejeitadas.} Os ângulos são reduzidos ao intervalo $]-\pi,\pi]$ e uma solução é rejeitada se:
\begin{itemize}
  \item alguma junta sair dos limites ($\pm165^\circ$ em J1 a J5 e $\pm175^\circ$ em J6, por omissão; em execução, os valores são lidos do robô);
  \item o cotovelo estiver demasiado dobrado, $|q_3|>140^\circ$, porque o antebraço fica quase encostado ao braço e a gravidade impede o robô de chegar ao ponto;
  \item o cotovelo, J4 ou J5 ficarem abaixo da base, isto é, se $z_c\le 0$, $z_{J4}=z_c+L_3\cos(q_2+q_3)\le 0$ ou $z_{J5}=z_{J4}+L_4\cos\varphi\le 0$.
\end{itemize}
O ponto é ainda inalcançável se $r\le L_5$ ou se $|\cos q_3|>1$ para todas as inclinações testadas.

\paragraph{Escolha da inclinação.} Para cada ponto começamos pela inclinação preferida, $\psi=90^\circ$. Se não houver solução, afastamo-nos de $5^\circ$ em $5^\circ$, primeiro para a base ($\psi>90^\circ$, útil para objetos perto da base) e depois para fora ($\psi<90^\circ$, para objetos longe), dentro do intervalo $[30^\circ,135^\circ]$. O limite inferior evita a garra a apontar quase na horizontal, que não serve para apanhar objetos, e o superior limita a inclinação para a base a $45^\circ$. Quando um candidato funciona, refinamo-lo por bisseção, para ficar com a inclinação mais próxima de $90^\circ$ que ainda tem solução, de modo a que a garra incline de forma contínua e não por saltos de $5^\circ$. Nos movimentos verticais usamos uma só inclinação, a pedida, para que uma descida inteira use a mesma solução, sem mudanças bruscas das juntas.

%---------------------------------------------------------------------
\subsection{Do referencial do mundo para o da base}

O CoppeliaSim dá as posições dos objetos no referencial do mundo, mas a cinemática trabalha no referencial da base, que está rodado em torno de $z$. Com $\mathbf{p}_b$ a posição da base no mundo e $\theta_b$ a sua rotação, a conversão é
\begin{equation*}
\mathbf{p}_{base}=R_z(-\theta_b)\,(\mathbf{p}_{mundo}-\mathbf{p}_b),
\end{equation*}
implementada em \texttt{aux.world\_to\_base}. No nosso cenário,
\begin{equation*}
\mathbf{p}_b=(-0{,}0497;\;0{,}0165;\;0{,}03)~\mathrm{m}, \qquad \theta_b=-90^\circ.
\end{equation*}
Não basta subtrair $\mathbf{p}_b$, porque os eixos da base também estão rodados; é preciso rodar o vetor de $-\theta_b$.

%---------------------------------------------------------------------
\subsection{Resultados para o nosso caso}

\paragraph{Exemplo 1: alvo com a garra na vertical.} Tomámos o ponto
\begin{equation*}
\mathbf{p}=(0;\;0{,}20;\;0{,}06)~\mathrm{m}
\end{equation*}
no referencial da base, com $\psi=90^\circ$, isto é, $\varphi=-90^\circ$. Os cálculos deram:
\begin{align*}
r &= 0{,}2000~\mathrm{m},\\
h &= \sqrt{0{,}2^2-0{,}0646^2}=0{,}18928~\mathrm{m},\\
z-L_1 &= 0{,}06-0{,}1316=-0{,}0716~\mathrm{m},\\
q_1 &= 90^\circ-\operatorname{atan2}(0{,}18928;\,0{,}0646)=90^\circ-71{,}16^\circ=18{,}84^\circ,\\
W_h &= 0{,}18928-0{,}0731=0{,}11618~\mathrm{m},\\
W_v &= -0{,}0716+0{,}1513=0{,}0797~\mathrm{m},\\
\cos q_3 &= -0{,}0733 \;\Rightarrow\; q_3=\pm94{,}21^\circ.
\end{align*}
Ficámos com o cotovelo para cima, $q_3=-94{,}21^\circ$, cujo cotovelo fica a $0{,}239~\mathrm{m}$ de altura, contra $0{,}116~\mathrm{m}$ na solução de cotovelo para baixo. Seguem-se $\alpha_1=77{,}26^\circ$, logo $q_2=-12{,}74^\circ$, e $q_4=-90^\circ-(-12{,}74^\circ)-(-94{,}21^\circ)=16{,}95^\circ$. O resultado foi
\begin{equation*}
\mathbf{q}=(18{,}84^\circ;\;-12{,}74^\circ;\;-94{,}21^\circ;\;16{,}95^\circ;\;0^\circ;\;0^\circ).
\end{equation*}
Aplicando a cinemática direta a estes ângulos, recuperámos o alvo,
\begin{equation*}
\mathbf{p}=(0;\;0{,}2000;\;0{,}0600)~\mathrm{m}, \qquad \mathbf{a}=(0,0,-1).
\end{equation*}
É esta a configuração usada como segundo exemplo da secção~\ref{sec:direta}, o que fecha o ciclo entre as duas cinemáticas.

\paragraph{Exemplo 2: alvo longe, que obriga a inclinar a garra.} Tomámos agora
\begin{equation*}
\mathbf{p}=(0{,}05;\;0{,}30;\;0{,}02)~\mathrm{m}.
\end{equation*}
Aqui, $r=0{,}3041~\mathrm{m}$, $h=0{,}2972~\mathrm{m}$ e $z-L_1=-0{,}1116~\mathrm{m}$. Com $\psi=90^\circ$ obtivemos $W_h=0{,}2241~\mathrm{m}$ e $W_v=0{,}0397~\mathrm{m}$, ou seja, o pulso ficaria a $0{,}2276~\mathrm{m}$ de J2, acima do máximo $L_2+L_3=0{,}2064~\mathrm{m}$, e $\cos q_3=1{,}434>1$: sem solução. Como o alvo está longe, inclinar a garra para a base afasta ainda mais o pulso, pelo que só funcionam inclinações $\psi<90^\circ$. A pesquisa parou em $\psi\approx82{,}27^\circ$, onde $\cos q_3\approx1$ e o braço fica quase esticado, o que mostra que o ponto está no limite do alcance. Os ângulos foram
\begin{equation*}
\mathbf{q}=(2{,}80^\circ;\;-81{,}96^\circ;\;-0{,}21^\circ;\;-0{,}10^\circ;\;0^\circ;\;0^\circ),
\end{equation*}
e a cinemática direta devolveu
\begin{align*}
\mathbf{p} &= (0{,}05;\;0{,}30;\;0{,}02)~\mathrm{m},\\
\mathbf{a} &= (-0{,}0066;\;0{,}1343;\;-0{,}9909),
\end{align*}
isto é, a garra aponta para baixo, ligeiramente para fora.

Nos dois casos, a cinemática direta devolveu o ponto pedido, o que valida a cinemática inversa.
